% =====================================================================
\section{The day and sleep}
\fig{fig_day.png}{What reaches its maximum and minimum, and when, over the day in rodents. Data from different studies, species and brain regions; for the bottom row the times are approximate.}
Almost every cell of the body runs a clock of its own with a period of about a day. Its mechanism is a negative-feedback loop with delay: two proteins switch on genes whose products, some hours later, suppress those same two proteins. The master clock sits in a small brain region above the crossing of the optic nerves; light through the eyes adjusts it to the alternation of day and night, and it synchronises the rest.

For our subject the day matters for two reasons. First, sleep and waking change neuronal activity and with it the influx of oxidants: the load on the defence is not constant, it passes through a cycle every day. Second, the oxidative defence itself oscillates with a daily period, and these oscillations have been found even in cells without a nucleus, that is, without genes.
\facts
\begin{itemize}
\item The clock's intrinsic period is preserved in constant darkness.\tU\ A daily rhythm of peroxiredoxin hyperoxidation is present in red blood cells, which have no nucleus.\tU
\item In the master clock of the mouse brain, the amount of hyperoxidised peroxiredoxin and the total amount of the enzyme both vary severalfold over the day; the hyperoxidised form leads the total amount by about 4 hours (figure reading; three animals).\tS
\item In slices of the master clock the tissue is reduced by day and oxidised at the beginning of the night; a change in the oxidative state alters neuronal excitability within about a minute.\tS
\item In PV cells of the rat frontal cortex a marker of oxidative DNA damage is maximal towards the end of the active phase and falls by 47\,\% by the end of the rest phase (fig.~8).\tP
\item Six hours of disrupted sleep raise this marker in rats by 25\,\%; in PV cells it is 27--50\,\% higher than in neighbouring cells.\tS\ After 4 hours without sleep the firing rate of cortical neurons is 34--64\,\% higher.\tS
\item In PV cells of the mouse frontal cortex 9\,\% of genes are read with a daily rhythm in males and 14\,\% in females (about 35\,\% in pyramidal cells); in males these are mainly genes for mitochondrial proteins. The intrinsic excitability of PV cells is the same by day and by night.\tP
\item In the mouse adrenal gland, heart and brown fat, mitochondrial hyperoxidised peroxiredoxin and the sulfiredoxin that repairs it oscillate with a daily period nearly in antiphase. Damage is maximal towards the end of the active phase, sulfiredoxin 8--12 hours later (figure reading; six points over the day). The total amount of peroxiredoxin itself is constant.\tS
\item Sleep and circadian disturbances occur in a substantial fraction of patients with schizophrenia.\tU
\end{itemize}
\begin{note}
Rodents are nocturnal: their active phase falls in darkness. Clock times cannot be carried over from mouse to human directly; what carries over is the phase relative to sleep and waking. The daily rhythm of hyperoxidised peroxiredoxin in the cortex and in PV cells has not been measured: the data are for the master clock and for organs outside the brain. In a whole-mouse-brain sample, mitochondrial hyperoxidised peroxiredoxin was not detected under ordinary conditions;\tS\ such a sample says nothing about a rare cell type.
\end{note}
\begin{ineq}
There are two distinct objects here. The first is a periodic external forcing of a system with two states: if the daily range of the load is comparable with the width of the window, the system approaches the boundary every day. The second is an intrinsic relaxation cycle ``damage $\to$ repair switched on $\to$ recovery $\to$ repair decays''; such a model has been published for cells outside the brain (del Olmo et al., 2019).\tS\ The two can be distinguished by the phase shift: a first-order element cannot lag its input by more than a quarter of the period, whereas the observed 8--12 hours out of 24 is between a third and a half.\tD
\end{ineq}
